\documentclass[11pt]{article}

\usepackage[margin=1.1in]{geometry}
\usepackage{amsmath,amssymb}
\usepackage{booktabs}
\usepackage{longtable}
\usepackage{array}
\usepackage{setspace}
\usepackage[hidelinks]{hyperref}

\setlength{\parskip}{0.35em}
\newcommand{\ind}[1]{\mathbf{1}\{#1\}}
\newcolumntype{L}[1]{>{\raggedright\arraybackslash}p{#1}}

\title{\textbf{Housing First or Treatment First?}\\[0.4em]
\large Model, calibration, validation, and mechanisms}
\author{}
\date{Draft --- \today}

\begin{document}
\maketitle

\begin{abstract}
\noindent This note documents the rebuilt model, its two-stage calibration, the
untargeted moments it is validated against, and the mechanism decomposition that
replaces the appendix's static model. It is written to be read as a draft of
Sections 2--5 of the revision. Numerical results from the calibration and the
mechanism ladder are marked where they are still pending; the specifications,
targets and sources are final.
\end{abstract}

\tableofcontents

%======================================================================
\section{The model}

\subsection{Environment}

Time is discrete and a period is one year. The economy is populated by
overlapping generations of households indexed by age $j=1,\dots,J$, corresponding
to real ages $26$ to $85$, with retirement at real age $66$ ($j_{R}=41$).
Households face age-, health- and behaviour-dependent mortality; the assets of
the deceased are redistributed to survivors as a uniform lump-sum bequest.

The individual state is
\[
  s=\bigl(j,\;a,\;y,\;e,\;d,\;m,\;\eta\bigr),
\]
with $a\ge0$ financial assets, $y$ labour productivity, $e$ employment and
unemployment-insurance status, $d\in\{0,1\}$ substance use disorder (SUD),
$m\in\{0,1\}$ serious mental illness (SMI), and $\eta\in\{0,1\}$ possession of a
housing voucher. Table~\ref{tab:state} gives the discretisation.

\begin{table}[t]
\centering
\caption{State space: $12{,}960$ points per age}
\label{tab:state}
\begin{tabular}{llrL{7.4cm}}
\toprule
State & Symbol & Points & Notes\\
\midrule
Age & $j$ & 60 & real age 26--85, retirement at 66\\
Financial assets & $a\ge0$ & 60 & cubic spacing, dense near zero; no borrowing\\
Labour productivity & $y$ & 9 & Rouwenhorst discretisation of an AR(1) in logs, $\rho=0.90$, $\sigma=0.20$\\
Employment & $e$ & 3 & employed; unemployed with UI; unemployed without UI\\
SUD status & $d$ & 2 & clean; substance use disorder\\
Serious mental illness & $m$ & 2 & absent; present\\
Housing voucher & $\eta$ & 2 & not held; held\\
\bottomrule
\end{tabular}
\end{table}

Two features of the state space are deliberate and are worth stating because
they are what make the model solvable at this size.

\paragraph{Housing entered with is not a state.} Survival, employment
transitions and the SUD recovery hazard all condition on the \emph{chosen}
housing $h'$, and renters face no adjustment cost, so the unit a household
arrives in is payoff-irrelevant. Carrying it would replicate the entire problem
$N_h$ times. Removing it leaves every aggregate unchanged and cuts the solve
time by a factor of $6.6$.

\paragraph{Three employment states, not four.} Unemployment insurance is drawn
on \emph{entry} into unemployment with probability $\pi_{UI}=0.30$ and lasts one
year. A first-year non-recipient and a household unemployed for two or more
years have identical income (zero) and identical continuation values, so those
two states merge \emph{exactly}. This is an equivalence, not an approximation,
and it answers the Associate Editor's question A2 about whether modelling UI
expiration is important: it is, but it costs one state rather than two.

\subsection{Preferences}

Period utility over consumption $c$, housing services $h$, drug consumption $v$
and hours $n$ is
\begin{equation}
u(c,h,v,n;d,m,h)=\frac{\bigl[c^{\alpha}h^{\beta}(v+b_{dm h})^{\gamma}\bigr]^{1-\sigma}}{1-\sigma}
+\omega_{\ell}\frac{(1-n)^{1-\sigma_{\ell}}}{1-\sigma_{\ell}}
+\aleph-\kappa_{x}\bigl(H+\omega_{x}H_{D}\bigr),
\label{eq:util}
\end{equation}
with
\[
\alpha=(1-\theta_h)(1-\theta_v),\qquad
\beta=\theta_h(1-\theta_v),\qquad
\gamma=\theta_v .
\]

\paragraph{Cobb--Douglas weights do not depend on SUD status.} This is the change
requested by Referee~2 (C5), Referee~3 (M2a) and the AE (A8). In the previous
specification the taste weight on drugs rose with $d$, so reducing SUD
mechanically raised the utility derived from $c$ and $h$, making
addiction-reduction and housing expansion complements ``even absent incentive
effects''. Holding $\alpha,\beta,\gamma$ fixed removes that channel entirely.

\paragraph{Only the Stone--Geary shifter varies.} $b_{dmh}$ depends on SUD
status, mental illness and \emph{chosen housing}:
\begin{equation}
b_{dmh}=b_d\cdot \underbrace{\rho_m^{\,\ind{m=1}}}_{\text{self-medication}}\cdot
\underbrace{\rho_H^{\,\ind{h=\text{homeless}}}}_{\text{housing access}} ,
\qquad \rho_m,\rho_H\in(0,1].
\label{eq:b}
\end{equation}
A lower $b$ raises the marginal utility of the first units of drug use, so poor
mental health and homelessness both push a household towards initiating. Three
remarks.

First, the AE's suggestion (A8) was to keep the Cobb--Douglas weights invariant
and vary $b$ instead. That is what \eqref{eq:b} does. A \emph{single} $b$ cannot
match both the participation moment and the drug-expenditure share: a $b$ large
enough to deter clean households pushes the unconstrained interior choice
$v^{*}$ below zero, which drives drug spending by households with SUD to the
corner. Two values, $b_0$ for clean and $b_1$ for SUD, separate the two moments
--- $b_0$ pins participation and $b_1$ keeps the choice interior so that
$\theta_v$ pins the expenditure share.

Second, $\rho_m<1$ is the self-medication channel Referee~1 asks about in
comment 1(a).

Third, and most consequentially, \textbf{$\rho_H$ is the only parameter in the
model that governs the elasticity of drug use with respect to housing access}.
Referee~2 calls this elasticity the crux on which all the welfare results rest;
Referee~3 and the AE (A13) say the same independently. Without $\rho_H$ the
moment is not merely untargeted, it is \emph{untargetable}: housing a household
raises its resources, drugs are a normal good, and the model would predict that
housing \emph{raises} use --- the opposite sign to Milby et al.\ (2005).
Section~\ref{sec:secondstage} calibrates $\rho_H$ to that trial.

\paragraph{The value of life.} $\aleph$ is the flow value of being alive. It is
what makes mortality risk deter drug use, and it therefore carries the paper's
central mechanism; Section~\ref{sec:mech} shows exactly how much work it does.

\paragraph{The externality.} $H$ is the aggregate homelessness rate and $H_D$
the share of the population both homeless and with SUD. Referee~2's wording is
specific --- ``a large homeless population, \emph{especially if many are
addicted to drugs}, creates a negative externality beyond the fiscal burden'' ---
so the disamenity is weighted by composition through $\omega_x$, not by the
headcount alone. Households take $H$ and $H_D$ as given, which is what makes
this a spillover rather than a private cost, but both are equilibrium objects,
so policy moves them and the movement feeds back through the value of life into
drug decisions. This is precisely the property the loss-of-consortium adjustment
$\kappa_{\text{cons}}$ lacked, and which Referee~1 objected to in comment 3(a).
The baseline sets $\kappa_x=0$; see Section~\ref{sec:extern}.

\subsection{Labour supply}

Hours are continuous and are pinned by the intratemporal first-order condition,
\begin{equation}
u_c(c)\cdot\frac{\partial R(n)}{\partial n}=\omega_{\ell}(1-n)^{-\sigma_{\ell}},
\label{eq:foc}
\end{equation}
where $R(n)$ denotes resources net of saving. Two features of \eqref{eq:foc}
matter economically, not just numerically.

\paragraph{The residual is not monotone.} The food-stamp floor is clawed back at
a rate of one and the housing voucher at the tenant contribution rate
$\kappa_r=0.30$, so $\partial R/\partial n$ \emph{jumps up} at the hours where
each transfer is exhausted, and the left-hand side of \eqref{eq:foc} can cross
the right-hand side more than once. This is the welfare trap: below the
food-stamp exhaustion point the return to an extra hour is close to zero, so
``work a little and keep the transfer'' and ``work properly'' are genuinely
distinct local optima and the household must compare them. The solution
algorithm splits $[0,1]$ at the kinks and solves each piece exactly rather than
bracketing once, which would silently select whichever optimum happened to be
bracketed.

\paragraph{Voucher eligibility is a step in the continuation value.} Income
eligibility is lost discretely at $50\%$ of median income. That does not kink
the budget, so \eqref{eq:foc} is continuous across it; it shifts the
continuation value, so the piecewise optimum is clamped to each side of the
eligibility threshold and both are evaluated.

Referee~3 (m1) wanted labour supply retained and its role made visible; the AE
(A3) questioned whether it was needed at all. Making it continuous and exact
resolves the tension in Referee~3's favour at a cost of a factor of three in
solve time.

\subsection{Drug use, addiction and treatment}

A household chooses whether to use ($v>0$) or abstain ($v=0$), and how much,
subject to the price $p_v$. A clean household that uses acquires SUD next period
with probability $\pi_d$, \emph{independent of how much it consumes}. There is
no light/heavy threshold and no consumption-dependent hazard; the previous
specification's $\bar v$ and the $6.48$ heavy-to-light ratio are gone. This is a
simplification, and it removes two of the parameters Referee~3 (M2b) asked to
have justified.

A household with SUD recovers with hazard
\begin{equation}
\pi_{r}(h,\tau,m)=\bar\pi_{r}\cdot \mu_H^{\,\ind{h=\text{homeless}}}\cdot
\mu_m^{\,\ind{m=1}}\cdot g^{\,\tau},
\end{equation}
where $\tau\in\{0,1\}$ is a discrete \textbf{treatment take-up} decision carrying
a resource cost, and $g>1$ is the treatment effect on the recovery hazard.
Referee~1's comment 1(b) and the AE's A10 are the same point and it is the most
damning gap in the submitted paper: the policy is called Treatment First and
there is no treatment in the model. There is now.

\subsection{Housing, vouchers, and the two policies}

Rent is a fixed price; there is no housing-market general equilibrium. This
follows the AE's A5 --- prices and rents cleared two markets whose PE and GE
outcomes were nearly identical. Homelessness is the lowest point of a four-point
housing grid with a calibrated service flow, replacing the seven-point grid with
four uncalibrated values that Referee~2 (D1) objected to.

Voucher rules are institutional and are not free parameters: income eligibility
below $50\%$ of median income, a $30\%$ tenant contribution, and a lottery
probability calibrated to HUD take-up among the eligible. Vouchers pay only
while renting.

The two policies differ in one rule:
\begin{center}
\begin{tabular}{lcc}
\toprule
& Housing First & Treatment First\\
\midrule
Voucher conditional on abstinence & no & yes\\
Detection probability & --- & $p_{\text{det}}\le1$\\
False-positive rate & --- & $p_{\text{fp}}\ge0$\\
\bottomrule
\end{tabular}
\end{center}
Referee~1 (4a, 4b) and the AE (A11) both note that perfect observation and
enforcement are not credible. $p_{\text{det}}$ and $p_{\text{fp}}$ are
experiment parameters reported as sensitivity, not calibrated, because that
is exactly the point.

\subsection{Serious mental illness}

$m$ is an exogenous persistent binary state representing \textbf{serious} mental
illness, not any mental illness. The distinction is not cosmetic: AMI covers
roughly $23\%$ of US adults and includes mild anxiety and depression, whereas
SMI is about $5.5\%$ and is defined by substantial functional impairment. It is
SMI that drives chronic homelessness and disability receipt, and a state that
halves productivity cannot sensibly cover a quarter of the population.

The specification follows \.{I}mrohoro\u{g}lu and Zhao (2022) structurally as
well as numerically. In their model labour-market shocks generate \emph{short}
homeless spells while health shocks generate \emph{chronic} homelessness, and
expanded means-tested disability income is the effective policy for the chronic
group. Homelessness driven purely by income risk is transitory; the persistent
tail requires a persistent state. SMI affects six objects: productivity
conditional on working, job finding and separation, mortality, the SUD recovery
hazard, the Stone--Geary shifter through $\rho_m$ in \eqref{eq:b}, and
eligibility for disability income when not employed. Onset is \emph{derived}
from the target prevalence and the recovery rate, so calibrating persistence
never disturbs prevalence.

The productivity penalty is also what generates \emph{employed} households poor
enough to be priced out of housing --- Referee~1's comment 2(c)(i), which an
unemployment effect alone cannot deliver --- and the persistence is what
generates the older homeless population of comment 2(c)(ii).

\subsection{Government and equilibrium}

A stationary equilibrium is a set of prices and aggregates such that households
optimise, the government budget balances, the drug market clears, bequests equal
the assets of the deceased, and the externality aggregates are consistent:

\begin{center}
\begin{tabular}{lL{9.6cm}}
\toprule
Object & Determined by\\
\midrule
$\tau_{\ell}$ & government budget balance, closed form given the distribution\\
$p_v$ & drug-market clearing, with production cost $C(V)=\frac{\psi_1}{\psi_2}V^{\psi_2}$, so the supply elasticity is $1/(\psi_2-1)$\\
Bequests & assets of the deceased, redistributed uniformly\\
$H,\;H_D$ & fixed point, since they enter period utility\\
\bottomrule
\end{tabular}
\end{center}

Making $p_v$ an equilibrium object answers Referee~1 (1c), Referee~3 (M2d) and
the AE (A5). Referee~3 worried that endogenising the price ``requires
suppressing this normalization''. It does not: $\psi_1$ is normalised so that
$p_v=1$ at the benchmark, which leaves the existing normalisation intact and
means only counterfactuals move the price.

Bequests are uniform rather than type-dependent, following the AE's A7.

%======================================================================
\section{Calibration}
\label{sec:calib}

The calibration has two stages, and they are different exercises. The first
stage matches a cross-section of the stationary equilibrium. The second stage
matches short-run \emph{responses} measured in experiments, which cannot be
computed from a comparison of stationary equilibria.

\subsection{First stage: the stationary equilibrium}

\subsubsection{Parameters set outside the model}

\begin{longtable}{lL{2.8cm}L{7.0cm}}
\caption{Externally set parameters}\label{tab:ext}\\
\toprule
Parameter & Value & Source\\
\midrule
\endfirsthead
\toprule
Parameter & Value & Source\\
\midrule
\endhead
$\sigma$ & 2.0 & standard\\
$\sigma_{\ell}$ & 3.0 & Frisch elasticity of 0.5, Chetty (2012)\\
$\xi$ (tax progressivity) & 0.85 & Heathcote, Storesletten and Violante (2017)\\
$r$ & 0.03 & Kaplan, Mitman and Violante (2020)\\
$\kappa_V$ (eligibility) & $0.50\times$ median & HUD programme rule\\
$\kappa_r$ (tenant share) & 0.30 & HUD programme rule\\
$\pi_{UI}$ & 0.30 & insured unemployment / total unemployment, FRED 2014--24\\
UI replacement rate & 0.50 & standard US\\
$\pi_{ue}$ (job finding) & 0.55 & Krueger, Cramer and Cho (2014)\\
$\kappa_{ue,H}$ & 0.80 & Desmond and Gershenson (2016)\\
$\kappa_{\phi,H}$ & 1.60 & Meyer et al.\ (2025): 60\% excess mortality when unhoused\\
$\kappa_{\phi,m}$ & 2.50 & standardised mortality ratio for SMI of 2--3\\
SMI prevalence & 5.5\% & NSDUH 2023, adults\\
SMI recovery hazard & 0.067/yr & \.{I}mrohoro\u{g}lu--Zhao's own moment: about half improve enough to work within ten years\\
$\chi_D$ & 0.212 & CEA (2017), roughly \$30{,}000 per person-year\\
$\chi_H$ & 0.252 & National Alliance to End Homelessness; reported as sensitivity, since R2 and R3 both judge it high\\
$\rho_{\text{pen},d}$ (SUD wage penalty) & 1.00 & no direct wage effect in the baseline; see below\\
\bottomrule
\end{longtable}

On the last row: the AE (A20) asks whether SUD has direct productivity effects.
The robust findings are on the \emph{employment} margin --- roughly $50\%$ more
unscheduled absence and $44\%$ higher turnover, which the model captures through
$\kappa_{ue,d}$ --- while estimates of the wage penalty conditional on working
are contested, with parts of the older literature finding positive effects for
casual use. The baseline therefore imposes no direct wage effect and says so,
rather than inventing a number.

\subsubsection{Internally calibrated parameters}

All parameters are calibrated jointly, but the parameter-to-moment map is close
to triangular: each parameter dominates one moment. Table~\ref{tab:int} states
the assignment. In the code the sign of $\partial M/\partial\theta$ is declared
explicitly for each pair, so a mistaken sign shows up immediately as a moment
walking away from its target rather than as a silent failure.

\begin{longtable}{lL{5.0cm}rL{4.0cm}}
\caption{Internally calibrated parameters and their moments}\label{tab:int}\\
\toprule
Parameter & Moment & Target & Source\\
\midrule
\endfirsthead
\toprule
Parameter & Moment & Target & Source\\
\midrule
\endhead
$b_0$ & drug participation & 5.13\% & NSDUH 2024, ages 26+\\
$\pi_d$ & SUD rate & 3.68\% & NSDUH 2024, ages 26+\\
$h_{\text{homeless}}$ & homelessness rate & 0.30\% & HUD point-in-time\\
$\pi_{eu}$ & unemployment rate & 4.90\% & BLS\\
$\kappa_{ue,d}$ & unemployment rate given SUD & 26.90\% & NSDUH\\
$\mu_H$ & drug use disorder share of the homeless & 21.70\% & Gutwinski et al.\ (2021)\\
$\pi_V$ & voucher recipients / income-eligible & 26.60\% & HUD\\
$\theta_v$ & drug spending / income given SUD & 44.00\% & existing moment\\
$\theta_h$ & aggregate housing cost / income & 21.20\% & SCF 2022, all households\\
$p_{\text{rent}}$ & share with housing cost above 50\% of income & 17.49\% & SCF 2022, all households\\
$\aleph$ & SUD income / average income & 68.70\% & NSDUH 2024\\
$\zeta_{DI}$ & SMI share of the homeless & 25.00\% & Gutwinski et al.\ (2021)\\
$\kappa_{ue,m}$ & unemployment rate given SMI & 7.10\% & NSDUH Table 6.5B, by Bayes\\
$\rho_{\text{pen},m}$ & homeless income / average & 19.10\% & Meyer et al., linked administrative data\\
$\beta$ & \textbf{wealth / GDP} & \textbf{3.00} & standard US capital--output ratio\\
$\underline c$ & food stamps / labour income & 0.42\% & SNAP outlays, FY2023\\
$\rho_m$ & SUD rate given SMI & 9.75\% & NSDUH 2023 detailed tables\\
$\omega_{\ell}$ & average hours worked & 40\% of time & Guvenen et al.\ (2023)\\
$b_1$ & keeps the drug choice interior, so $\theta_v$ identifies the expenditure share & --- & ---\\
\bottomrule
\end{longtable}

\paragraph{The discount factor and wealth/GDP.} $\beta$ is calibrated so that
aggregate wealth is three times GDP, with GDP measured as \emph{all factor
income} --- labour earnings plus the return on the capital stock --- and
transfers excluded, since transfers are redistribution rather than output. Three
is the standard US ratio of the private capital stock to annual output and is
the moment a life-cycle discount factor is normally pinned by.

This is the right target here rather than the SCF's net-worth-to-household-income
ratio of roughly $7.3$, whose numerator is mostly housing equity that a model
without homeownership cannot hold and whose denominator is a household income
concept rather than output. One implication should be stated in the paper rather
than left for a reader to find: with $A=3\cdot\text{GDP}$ and
$\text{GDP}=\text{labour income}+rA$, capital income is $9\%$ of GDP, which is
low against a one-third capital share. That follows from holding $r$ fixed at
$3\%$ in a model with no production function.

\paragraph{Three targets that were wrong in the submitted paper.} These are
worth flagging because two referees used them to distrust the quantification.

\begin{enumerate}
\item \emph{SUD among the homeless.} The paper targets $37\%$. Gutwinski et
al.\ (2021) pool 39 studies of homeless people in high-income countries and
report $36.7\%$ \emph{alcohol} use disorder and $21.7\%$ \emph{drug} use
disorder. The model is about drug use, so the target is $21.7\%$. This answers
the AE's A16.

\item \emph{Income of drug users.} Referee~1 (2c-i) attacked the model's $41\%$
against ``roughly $74\%$'' in the data. The $74\%$ is computed over people with
a use \emph{disorder}, not people who use --- a different population. For ages
26+, NSDUH gives $95.1\%$ for users and $68.7\%$ for moderate or severe use
disorder. Targeting the disorder ratio and \emph{validating} on the use ratio is
the sharper test, and it makes the economics explicit: use is close to
income-neutral, and it is progression to disorder that concentrates at the
bottom of the distribution.

\item \emph{Housing burden.} Everyone rents in this model, so a renters-only
target would import the selection into who rents. The target is therefore the
aggregate housing-cost-to-income ratio over \emph{all} households, with owner
housing valued at imputed rent, and the full distribution is used for validation
(Section~\ref{sec:valid}).
\end{enumerate}

\subsection{Second stage: short-run experimental moments}
\label{sec:secondstage}

Every report converges on the same objection. Referee~2 (C1): ``all of the
welfare results rest on whether the model has an accurate elasticity of drug use
to housing access''. Referee~3 and the AE (A13) say the same independently.
Referee~1 (2d, 2e) asks for policy-relevant elasticities from the RCT literature
and notes that the price-elasticity moment as written is not replicable.

These are short-run objects. A randomised trial hands a treatment to a selected
group and measures them months later; it does not move aggregate prices and it
does not wait for a new stationary distribution. The second stage therefore
holds prices at their benchmark values, takes the benchmark distribution as the
initial condition, restricts attention to the trial's eligible population,
applies each arm's rules, re-solves the household problem, and pushes the cohort
forward.

\subsubsection{The housing-access elasticity}

Milby et al.\ (2005) randomised 196 cocaine-dependent homeless adults into three
arms, all receiving identical day treatment:

\begin{center}
\begin{tabular}{llrl}
\toprule
Arm & Meaning here & $n$ & Weeks abstinent (of 24)\\
\midrule
NH   & no assistance & 66 & 4.48\\
NACH & housing, unconditional --- \emph{Housing First} & 67 & 6.60\\
ACH  & housing conditional on abstinence --- \emph{Treatment First} & 63 & 8.32\\
\bottomrule
\end{tabular}
\end{center}

The model counterparts are the same three arms applied to the model's homeless
population with SUD, and the moment targeted is the \emph{ratio}
$\text{NACH}/\text{NH}=1.473$. Ratios rather than levels: ``mean consecutive
weeks abstinent out of 24'' is not the same object as ``abstains this model
year'', but the treatment effect is comparable, and it is the treatment effect
the referees are asking the model to reproduce. This moment identifies $\rho_H$
in \eqref{eq:b}.

\paragraph{The conditionality effect is a validation check, not a target.} The
abstract of Milby et al.\ reports that the abstinence-contingent arm beat the
unconditional arm, but that holds only \emph{after controlling for treatment
attendance}. In the raw randomised comparison NACH versus ACH gives $P=.091$,
and the GEE models state that the two housed groups ``did not differ from one
another''. Calibrating to a point estimate the source paper reports as
insignificant would be exactly the error Referee~1 attacked
$\kappa_{\text{cons}}$ for. The model's $\text{ACH}/\text{NACH}$ ratio is
therefore reported against $1.261$ as a sign check with the confidence band
shown.

\paragraph{The extensive/intensive contrast is the sharper test.} Milby's
counterfactual is literally no housing, and the effect is large. The Housing
First RCTs compare Housing First to ``usual services'', which already include
shelter or transitional housing, and the systematic review finds ``no clear
differences \dots\ on outcomes related to mental health, quality of life and
substance use''. These are not in conflict. The elasticity of drug use with
respect to housing access is \textbf{large at the extensive margin} (nothing to
something) and \textbf{near zero at the intensive margin} (some housing to
permanent housing). A model that reproduces both is far better identified than
one that matches either number alone, and reproducing both is the test this
paper should set itself.

\subsubsection{The drug-price participation elasticity}

Referee~1 (2e) says this moment is not replicable as written. It now has an
unambiguous model counterpart: perturb $p_v$, hold everything else at the
benchmark, re-solve the household problem, and read the arc elasticity of the
participation rate. It is a partial-equilibrium perturbation deliberately, since
the empirical estimates come from price variation rather than from a comparison
of steady states. The target is $-0.63$ and it identifies $\kappa_{\phi,d}$, the
SUD mortality multiplier, whose admissible range is the $2.5$--$9.8$
standardised-mortality-ratio band of Degenhardt et al.\ (2011). Referee~1 also
asks for robustness across the $-0.36$ to $-0.9$ range, which is a re-run of the
same bisection at the two endpoints.

\subsubsection{The treatment effect and the value of a statistical life}

The treatment effect on the recovery hazard, $g$, is the parameter through which
the entire Treatment First result runs, so it deserves the same discipline
Referee~1 forced on $\kappa_{\text{cons}}$. The submitted value of $3.0$ ---
tripling the recovery hazard --- is not supported. Remission from SUD overall
runs $35$--$54\%$ at a mean 17-year follow-up and remission following specialty
treatment about $48\%$, but those are different populations with different
selection, not treatment effects. The only clean randomised estimate is an odds
ratio near $1.22$ for contingency management, which implies $g\approx1.35$. The
paper should report the Treatment First result across a range and identify the
value of $g$ at which the ranking would change.

The model-implied value of a statistical life is the marginal rate of
substitution between consumption and the survival probability. It is currently
\emph{reported} rather than targeted, because $\aleph$ carries a first-stage
moment; the reference value is \$7.5m (Viscusi and Aldy 2003; EPA 2010). If the
model-implied value is far from that, the right response is to move $\aleph$ to
the VSL target and reassign the SUD income ratio to $\rho_{\text{pen},d}$.

%======================================================================
\section{Validation}
\label{sec:valid}

None of the moments in this section is targeted. Referee~1's comment 2(a) ---
that model fit is never reported --- and the AE's A12 are answered by reporting
all of them, including the misses, and by framing the misses honestly rather
than favourably (comment 2b).

\subsection{Distributional moments}

\begin{longtable}{L{6.0cm}L{4.0cm}L{4.4cm}}
\caption{Untargeted distributional moments}\label{tab:valid1}\\
\toprule
Moment & Data & Source\\
\midrule
\endfirsthead
\toprule
Moment & Data & Source\\
\midrule
\endhead
Distribution of housing cost / income, 12 thresholds & 95.4, 89.3, 77.2, 63.6, 50.9, 40.6, 32.5, 26.0, 17.5, 12.6, 7.1, 4.6\% above 5--100\% & SCF 2022, ages 26--85, owner housing at imputed rent\\
Income of drug \emph{users} relative to average & 95.1\% & NSDUH 2024, ages 26+\\
Share of the homeless who are employed & 40--52\% & Referee~1, comment 2(c)(i)\\
Age distribution of the homeless & 48.1\% aged 26--44; 27.1\% aged 55+ & Referee~1, comment 2(c)(ii)\\
Drug use by age (U-shaped) & NSDUH & microdata\\
Homelessness by age & California point-in-time counts & administrative\\
Wealth by age & SCF 2022 & microdata\\
\bottomrule
\end{longtable}

\subsection{Policy responses}

\begin{longtable}{L{5.6cm}L{4.0cm}L{4.8cm}}
\caption{Untargeted policy responses}\label{tab:valid2}\\
\toprule
Response & Data & Source\\
\midrule
\endfirsthead
\toprule
Response & Data & Source\\
\midrule
\endhead
Housing First effect on 3-year mortality & $-4.6$pp (s.e.\ 1.7) & Costantini (2025), $\approx$300,000 VA veterans, IV\\
Treatment First effect on long-run health & $\approx0$ & Costantini (2025)\\
Housing stability under Housing First & sustained & HF RCT literature\\
Housing stability under Treatment First & declining & HF RCT literature\\
Property values: permanent supportive housing versus emergency shelter & PSH raises, shelter lowers & shelter-siting hedonic literature\\
\bottomrule
\end{longtable}

\paragraph{The Costantini result is the paper's largest risk and belongs in the
introduction.} The best available quasi-experimental evidence says the mortality
benefit runs through \emph{housing}, and that conditioning on treatment adds
nothing detectable. The submitted paper concludes that universal Treatment First
dominates. The defensible position after the revision is the \emph{generosity}
result --- expanding coverage dominates, and conditionality is second-order and
uncertain --- which is consistent with all of the evidence above and is still a
paper. Getting there requires reporting welfare and mortality as two separate
dimensions, which is the AE's A9, rather than collapsing them into one ranking.

%======================================================================
\section{Mechanisms: shutting channels and recalibrating $b$}
\label{sec:mech}

\subsection{Why this replaces the static appendix}

Referee~1's comment 1(d) does not merely ask for intuition, it \emph{specifies
the design}: show that the mortality channel is the quantitatively important
one by shutting channels off. Comment 1(e) says the current appendix is
disorganised, static, and overreaches. Referee~2 (C4) wants the drug-use policy
function shown by income and wealth for households with and without SUD.
Referee~3 (M2b) asks what drug-consumption patterns the Stone--Geary
specification actually generates and where heterogeneity in $v$ comes from;
(M2c) asks the sharpest theoretical question in any of the reports --- since
drug use is rational and forward-looking, where do the externalities justifying
intervention arise? The AE (A1) asks for a stylised section because the number
of moving parts obscures the economics.

One exercise answers all five, and it is run on the \emph{calibrated} model
rather than on a separate static toy. That matters: a static model with no
assets cannot speak to the question Referee~2 raised about wealth, and a toy
calibrated to nothing cannot establish that a channel is quantitatively
important.

\subsection{The design}

Shut one channel off at a time and, in each arm, \textbf{recalibrate $b_0$ so
that drug participation returns to its benchmark value}. The recalibration is
what makes the comparison informative. Without it, removing the mortality
penalty simply raises drug use everywhere, and the level change swamps any
compositional change. Holding participation fixed separates \emph{sorting} from
\emph{level} and yields two distinct statistics per arm:

\begin{center}
\begin{tabular}{ll}
\toprule
Statistic & What it measures\\
\midrule
How far $b_0$ had to move & how much deterrence the channel was providing\\
The income gradient of use & \emph{whom} the channel deterred\\
\bottomrule
\end{tabular}
\end{center}

The arms are run in partial equilibrium at benchmark prices. The object of
interest is the behavioural mechanism, not equilibrium feedback, and holding
prices fixed isolates it --- as well as being roughly fifty times cheaper.

\subsection{The arms}

\begin{enumerate}
\item \textbf{Benchmark.} All channels on.
\item \textbf{No mortality penalty} ($\kappa_{\phi,d}=1$). SUD no longer raises
the death hazard. This is Referee~1's channel of interest and the one the paper
claims is central.
\item \textbf{No employment penalty} ($\kappa_{ue,d}=1$). SUD no longer worsens
job finding or separation.
\item \textbf{No homelessness feedback} ($\kappa_{\phi,H}=\kappa_{ue,H}=\mu_H=1$).
Being homeless no longer raises mortality, worsens employment, or lowers the
recovery hazard. This arm isolates the loop that makes homelessness and SUD
mutually reinforcing.
\item \textbf{All private costs off at once.} Mortality, employment and the
homelessness loop are shut simultaneously. Whatever deterrence survives is
\emph{resource diversion}: the budget constraint alone. Diversion cannot be
switched off directly --- drugs are priced inside a Cobb--Douglas bundle, so
making them free sends $v$ to the corner --- but it is measurable as the
residual, which is what this arm reports.
\end{enumerate}

\subsection{What is reported}

For each arm, Table~\ref{tab:mech} reports the recalibrated $b_0$ relative to
benchmark, and then the composition of drug users: the use rate within each
income quintile, mean user income relative to average (split into experimenters
and households with SUD), and the shares of users who are unemployed and
homeless, together with their mean age and assets.

\begin{table}[t]
\centering
\caption{Mechanism decomposition (structure; numbers pending calibration)}
\label{tab:mech}
\begin{tabular}{lccccc}
\toprule
& Benchmark & No mortality & No employment & No homeless & All private\\
& & penalty & penalty & feedback & costs off\\
\midrule
$b_0$ relative to benchmark & 1.00 & & & &\\
Participation (\%, held fixed) & & & & &\\
SUD rate (\%) & & & & &\\
Homelessness (\%) & & & & &\\
\addlinespace
Use rate, income Q1 (\%) & & & & &\\
Use rate, income Q5 (\%) & & & & &\\
\addlinespace
User income / average & & & & &\\
\quad experimenters only & & & & &\\
\quad SUD only & & & & &\\
Users unemployed (\%) & & & & &\\
Users homeless (\%) & & & & &\\
Mean age of users & & & & &\\
Mean assets of users & & & & &\\
\bottomrule
\end{tabular}
\end{table}

\subsection{What the exercise is expected to show, and how it could fail}

The paper's thesis is that the private cost of SUD runs through mortality,
employment and resource diversion, and that this is why forward-looking
households with more to lose --- higher income, more assets, better health ---
abstain. Fixing preferences across SUD status (Section~\ref{sec:calib}) makes
this a commitment rather than an assumption: with $\alpha,\beta,\gamma$
invariant, becoming addicted no longer relaxes any constraint, so \emph{all} of
the deterrence must come from those three channels.

The exercise therefore has a sharp prediction and a sharp failure mode. If
removing the mortality penalty requires only a small change in $b_0$ and leaves
the income gradient of use nearly flat, the paper's central mechanism is not
carrying the weight the text claims, and the honest response is to say so and
re-centre the paper on the generosity result. If instead $b_0$ must fall
substantially and the income gradient inverts, the mechanism is established
numerically on the calibrated model --- which is exactly what Referee~1 asked
for and what the static appendix could not deliver.

\subsection{Where the externalities come from}
\label{sec:extern}

Referee~3's M2c asks where the externalities justifying intervention arise, given
that drug use is rational and forward-looking. They are of two kinds and the
shutdown ladder quantifies both.

\emph{Fiscal.} A household that becomes homeless or acquires SUD imposes
$\chi_H$ and $\chi_D$ on the government budget and reduces the tax base through
lower earnings. None of this is internalised, and all of it is measured directly
by the model.

\emph{Non-fiscal.} The disamenity term $\kappa_x(H+\omega_x H_D)$ in
\eqref{eq:util}. Referee~2 (C2) asked for this and said it was easy to
introduce; it is, but it has no credible point estimate. Hedonic shelter-siting
studies measure a \emph{local, avoidable} disamenity --- what people pay to live
further away --- and this model has no geography, so nobody can avoid anything
and a capitalisation coefficient does not map onto an aggregate public bad.
Corbae, Glover and Nattinger's externality is about neighbourhood housing
\emph{quality} rather than homelessness, and their identification comes from a
rent-decontrol experiment, so what transfers is their identification logic and
not their functional form.

The recommendation is therefore to set $\kappa_x=0$ in the baseline, report a
bounded sensitivity, and use the qualitative finding from that literature ---
that permanent supportive housing \emph{raises} nearby property values while
temporary emergency shelter lowers them --- as the substantive point. Permanent
supportive housing is Housing First-like and transitional shelter is Treatment
First-like, so the externality plausibly differs in \emph{sign} between the two
policies being compared. That is a stronger argument than a made-up parameter,
and note that $\kappa_x=0$ is a defensible baseline in a way that
$\kappa_{\text{cons}}=0$ was not: it was the value at which the paper's
conclusion reversed.

%======================================================================
\section{Status}

\paragraph{Done.} Household problem with continuous labour supply from the
first-order condition; stationary distribution; aggregates; the three
equilibrium fixed points; the calibration framework with all nineteen
parameter--moment assignments; the second-stage experiment machinery; the
mechanism ladder.

\paragraph{Running.} The first-stage calibration. Every number the model
currently prints is a placeholder and none should be read as a result.

\paragraph{Not yet built.} The transition-path solver. Referee~1's comment 3(d)
calls the several-period partial-equilibrium exercise ``a half-measure'' and
asks for proper transition paths; with no housing market to clear, a path
iterates on $\tau_{\ell}$ and $p_v$ alone, which is far cheaper than in the
original model. This is the critical path for the revision, because the
second-stage moments of Section~\ref{sec:secondstage} are short-run objects.

\paragraph{Open decisions.} Whether to add a parental warm-glow term to answer
Referee~2's C3 and the AE's A6 together, or simply to delete the intergenerational
language; and whether to retain the maximum experimentation age, which Referee~1
correctly notes is in tension with the U-shaped age profile of drug use.

\end{document}
